\begin{enumerate}
	\item le jour de ses quatorze ans, Juliette mesure environ \np[~cm]{160} (c'est le point B du graphique).

	\item L'âge de 12 ans correspond au point A. On constate que la pente de la tangente en ce point est plus forte que celle de la tangente en B : donc Juliette grandissait plus vite à 12 ans qu'à 14 ans.

	\item Le point C correspond à l'âge de 17 ans, et en ce point la tangente est manifestement horizontale : cela signifie que Juliette ne grandit plus et que sa croissance est terminée.

	\item Puisque la croissance de Juliette est terminée à 17 ans, elle gardera comme taille adulte la taille de ses 17 ans : environ \np[~cm]{163}.
\end{enumerate}
